% GENERATED FILE. Edit canonical lessons, then run npm run generate:books.
% canonical-source-hash: fnv1a64:db350100bdfb3b4e
% canonical-lessons: FR-C446-cordonnier, FR-C446-ressemeler, FR-C446-cirer, FR-C446-serrurier, FR-C446-recoudre

\chapter{Chez le cordonnier}
\label{ch:fr-chez-le-cordonnier}
\hlchaptermodality{446}

\begin{chapteropening}
\textbf{By the end of this chapter:} \emph{I can get things mended: the cobbler, resoling, polishing, the locksmith and sewing back.}

Complete the last lesson of chapter 446: I can get things mended: the cobbler, resoling, polishing, the locksmith and sewing back.
\end{chapteropening}

\section[le cordonnier]{le cordonnier — the cobbler, the shoe mender}

\label{lesson:FR-C446-cordonnier}



\begin{quote}
Before the new one: say the French for the haberdasher’s, then the French for the wine merchant.
\end{quote}

\subsection*{You'll want to know: le cordonnier}
\textbf{le cordonnier} — \textquotedblleft{}the cobbler, the shoe mender\textquotedblright{}.

\begin{quote}
Le cordonnier a réparé mes bottes en une journée.
\end{quote}

\begin{quote}
\emph{The cobbler mended my boots in a day.}
\end{quote}

\emph{Les cordonniers sont toujours les plus mal chaussés} — the cobbler's children go barefoot.

\begin{cousinweb}
\emph{cordonnier} is an altered form of an older word for a worker in fine leather, named after a city famous for it. EtymDB lists it without a source.
\end{cousinweb}

\begin{grammarlens}[title={What you've built}]
One more word for telling, explaining and describing at B1.
\end{grammarlens}

\subsection*{Your turn}
\begin{itemize}
\raggedright
  \item \emph{Say it:} \emph{le cordonnier}
  \item \emph{Say it:} \emph{le cordonnier}, once more
  \item \emph{Say it:} a sentence of your own with \emph{le cordonnier}
  \item \emph{Recall:} say the French for the haberdasher’s, then the French for the wine merchant, then say \emph{le cordonnier} again
\end{itemize}

\begin{tcolorbox}[breakable,colback=teal!4,colframe=teal!35!black,title={Before you move on}]
What does le cordonnier mean? (\textquotedblleft{}The cobbler, the shoe mender\textquotedblright{}.) Say it once more.
\end{tcolorbox}
\section[ressemeler]{ressemeler — to resole}

\label{lesson:FR-C446-ressemeler}



\begin{quote}
Before the new one: say the French for the wine merchant, then the French for the cobbler.
\end{quote}

\subsection*{You'll want to know: ressemeler}
\textbf{ressemeler} — \textquotedblleft{}to resole\textquotedblright{}.

\begin{quote}
Il fait ressemeler ses chaussures chaque hiver.
\end{quote}

\begin{quote}
\emph{He has his shoes resoled every winter.}
\end{quote}

\emph{La semelle} — the sole.

\begin{cousinweb}
\emph{ressemeler} is built from \emph{re-} 'again' and \emph{semelle} 'sole' (EtymDB lists the verb). The origin of \emph{semelle} is not settled.
\end{cousinweb}

\begin{grammarlens}[title={What you've built}]
One more word for telling, explaining and describing at B1.
\end{grammarlens}

\subsection*{Your turn}
\begin{itemize}
\raggedright
  \item \emph{Say it:} \emph{ressemeler}
  \item \emph{Say it:} \emph{ressemeler}, once more
  \item \emph{Say it:} a sentence of your own with \emph{ressemeler}
  \item \emph{Recall:} say the French for the wine merchant, then the French for the cobbler, then say \emph{ressemeler} again
\end{itemize}

\begin{tcolorbox}[breakable,colback=teal!4,colframe=teal!35!black,title={Before you move on}]
What does ressemeler mean? (\textquotedblleft{}To resole\textquotedblright{}.) Say it once more.
\end{tcolorbox}
\section[cirer]{cirer — to polish, to wax}

\label{lesson:FR-C446-cirer}



\begin{quote}
Before the new one: say the French for the cobbler, then the French for to resole.
\end{quote}

\subsection*{You'll want to know: cirer}
\textbf{cirer} — \textquotedblleft{}to polish, to wax\textquotedblright{}.

\begin{quote}
Le dimanche, mon père cire ses chaussures.
\end{quote}

\begin{quote}
\emph{On Sundays my father polishes his shoes.}
\end{quote}

\emph{La cire} — wax; \emph{le cirage} — shoe polish.

\begin{cousinweb}
\emph{cirer} is built on \emph{cire} 'wax', which is inherited from Latin \emph{cēra} (Lewis \& Short).
\end{cousinweb}

\begin{grammarlens}[title={What you've built}]
One more word for telling, explaining and describing at B1.
\end{grammarlens}

\subsection*{Your turn}
\begin{itemize}
\raggedright
  \item \emph{Say it:} \emph{cirer}
  \item \emph{Say it:} \emph{cirer}, once more
  \item \emph{Say it:} a sentence of your own with \emph{cirer}
  \item \emph{Recall:} say the French for the cobbler, then the French for to resole, then say \emph{cirer} again
\end{itemize}

\begin{tcolorbox}[breakable,colback=teal!4,colframe=teal!35!black,title={Before you move on}]
What does cirer mean? (\textquotedblleft{}To polish, to wax\textquotedblright{}.) Say it once more.
\end{tcolorbox}
\section[le serrurier]{le serrurier — the locksmith}

\label{lesson:FR-C446-serrurier}



\begin{quote}
Before the new one: say the French for to resole, then the French for to polish.
\end{quote}

\subsection*{You'll want to know: le serrurier}
\textbf{le serrurier} — \textquotedblleft{}the locksmith\textquotedblright{}.

\begin{quote}
Elle a perdu ses clés et elle a appelé un serrurier.
\end{quote}

\begin{quote}
\emph{She lost her keys and called a locksmith.}
\end{quote}

\emph{La serrure} — the lock.

\begin{cousinweb}
\emph{serrurier} is built on \emph{serrure} 'lock', from \emph{serrer} 'to close tight', which goes back to Vulgar Latin \emph{serrāre}, built on Latin \emph{sera} 'bar for a door' (Lewis \& Short). EtymDB does not list it.
\end{cousinweb}

\begin{grammarlens}[title={What you've built}]
One more word for telling, explaining and describing at B1.
\end{grammarlens}

\subsection*{Your turn}
\begin{itemize}
\raggedright
  \item \emph{Say it:} \emph{le serrurier}
  \item \emph{Say it:} \emph{le serrurier}, once more
  \item \emph{Say it:} a sentence of your own with \emph{le serrurier}
  \item \emph{Recall:} say the French for to resole, then the French for to polish, then say \emph{le serrurier} again
\end{itemize}

\begin{tcolorbox}[breakable,colback=teal!4,colframe=teal!35!black,title={Before you move on}]
What does le serrurier mean? (\textquotedblleft{}The locksmith\textquotedblright{}.) Say it once more.
\end{tcolorbox}
\section[recoudre]{recoudre — to sew back on, to stitch up}

\label{lesson:FR-C446-recoudre}



\begin{quote}
Before the new one: say the French for to polish, then the French for the locksmith.
\end{quote}

\subsection*{You'll want to know: recoudre}
\textbf{recoudre} — \textquotedblleft{}to sew back on, to stitch up\textquotedblright{}.

\begin{quote}
La dame a recousu le bouton de ma veste.
\end{quote}

\begin{quote}
\emph{The lady sewed the button back on my jacket.}
\end{quote}

Irregular: \emph{je recouds, il a recousu}.

\begin{cousinweb}
\emph{recoudre} is built from \emph{re-} 'again' and \emph{coudre} 'to sew', which goes back to Latin \emph{consuere}, from \emph{suere} 'to sew' (Lewis \& Short). EtymDB does not list it.
\end{cousinweb}

\begin{grammarlens}[title={What you've built}]
One more word for telling, explaining and describing at B1.
\end{grammarlens}

\subsection*{Your turn}
\begin{itemize}
\raggedright
  \item \emph{Say it:} \emph{recoudre}
  \item \emph{Say it:} \emph{recoudre}, once more
  \item \emph{Say it:} a sentence of your own with \emph{recoudre}
  \item \emph{Recall:} say the French for to polish, then the French for the locksmith, then say \emph{recoudre} again
\end{itemize}

\begin{tcolorbox}[breakable,colback=teal!4,colframe=teal!35!black,title={Before you move on}]
What does recoudre mean? (\textquotedblleft{}To sew back on, to stitch up\textquotedblright{}.) Say it once more.
\end{tcolorbox}
